\subsubsection{Carry Handling in Microcode}
\label{sec:c25519-carries}

The implementation depends on how Goldmont microcode represents and consumes arithmetic carry state.
This mechanism differs from the conventional x86 interface and motivates the arithmetic organization used by both Curve25519 implementations.

An arithmetic $\mu$op such as \code{ADD} associates conditional state with its destination register, provided that destination is one of the sixteen internal \code{TMP} registers.
For an addition, this state records whether the operation produced a carry.
An architectural destination receives the correct arithmetic result but carries no state that the patch can query, in this domain or in \code{RFLAGS}.
Every carry producing addition in our patches therefore writes a \code{TMP}, and sums are copied out to architectural registers afterwards.
The state belongs to the register rather than to the addition, and it survives across triads until the next $\mu$op writes that register.
Any such write replaces it: a shift installs its own shifted out bit, and a multiply replaces the state of both registers it writes.
Because the state is held per register, several carries can be outstanding at the same time and can be recovered in any order.

The patch can consume this state in two ways.

The first path uses \code{SETCC\_CONDB}.
This operation reads the conditional state associated with a previous arithmetic destination and materializes the carry as an ordinary integer equal to zero or one.
The requirement to use an internal register applies only to the register being read; \code{SETCC\_CONDB} may write its result directly to an architectural register, as a full width zero extending write.
The patch can then add this value to any accumulator with a subsequent \code{ADD}.
The carry therefore becomes an explicit register value, at a cost of three $\mu$ops per propagated carry.
Different additions can materialize their carries into different registers, so several accumulation chains can progress without sharing \code{RFLAGS}.

The second path reconstructs the conventional x86 carry mechanism.
A microcode \code{ADC} consumes the carry flag in \code{RFLAGS} rather than the conditional state associated with an arithmetic destination.
That flag is captured from \code{RFLAGS} when the hook fires and is then frozen for the lifetime of the patch: no \code{ADD}, and no earlier \code{ADC}, updates it, so two \code{ADC}s placed back to back observe the same value.
To use a carry produced by a preceding \code{ADD}, the patch first publishes it with the two operand form \code{GENARITHFLAGS\_RR}, naming the \code{TMP} whose carry is wanted as the first operand.
\code{GENARITHFLAGS} transfers the associated condition into \code{RFLAGS}, after which \code{ADC} adds its two operands together with the architectural carry flag.
The operand form is load bearing.
The one operand \code{GENARITHFLAGS\_R} does not bridge; it writes a second architectural flag image that \code{READAFLAGS} observes and \code{ADC} ignores.
The register named must also be internal, since an architectural destination has no state to publish, so every accumulator in a chain has to be staged through a \code{TMP} and copied back at the end.
The destination of an \code{ADC} acquires conditional state of its own, so a fresh \code{GENARITHFLAGS\_RR} extends the chain; we verified chains of this shape to thirty two limbs, at roughly eight cycles per limb.
Each carry therefore costs two $\mu$ops, and the bridge and the following \code{ADC} can be packed into the same triad as the addition that feeds them.
We found no other write path into the register \code{ADC} reads, having probed \code{MOVEINSERTFLGS}, \code{MOVEMERGEFLGS} and all six \code{MOVETOCREG} variants directed at \code{CORE\_CR\_EFLAGS}.

Table~\ref{tab:c25519-carrypaths} summarizes the two paths.
Our $5\times51$ implementation uses the first.
The saturated $4\times64$ control uses the second.

\begin{table}[htbp]
\centering
\caption{Two ways to consume carry state in Goldmont microcode.
The main $5\times51$ implementation materialises carries as integers and accumulates them explicitly.
The saturated $4\times64$ control instead publishes carry state to \code{RFLAGS} and consumes it with \code{ADC}.}
\label{tab:c25519-carrypaths}
\footnotesize
\setlength{\tabcolsep}{5pt}
\begin{tabularx}{\linewidth}{@{}>{\raggedright\arraybackslash}p{.22\linewidth}>{\raggedright\arraybackslash}X>{\raggedright\arraybackslash}X@{}}
\toprule
Operation path & Effect on carry state & Use in our implementations \\
\midrule
\code{ADD} $\rightarrow$ \code{SETCC\_CONDB} $\rightarrow$ \code{ADD}
&
An arithmetic $\mu$op records conditional state with its internal destination register, where it stays until the next write to that register.
\code{SETCC\_CONDB} reads that destination associated state and materialises the carry as an integer equal to zero or one, in an internal or an architectural register.
A later \code{ADD} accumulates this value in a separate register.
Three $\mu$ops per carry, and the architectural flags are never touched.
&
Used in the main $5\times51$ multiplication and squaring patches, twenty five and fifteen times respectively, with no \code{ADC}.
This keeps the low sum, high sum, and carry count as separate dependency chains. \\
\midrule
\code{ADD} $\rightarrow$ \code{GENARITHFLAGS\_RR} $\rightarrow$ \code{ADC}
&
The arithmetic $\mu$op again records conditional state with its internal destination.
\code{GENARITHFLAGS\_RR}, naming that destination as its first operand, transfers the state into \code{RFLAGS}, which arithmetic $\mu$ops otherwise leave at the value captured when the hook fired.
A following \code{ADC} then consumes the architectural carry flag in the conventional way, and its own destination carries state that extends the chain.
Two $\mu$ops per carry, and every destination in the chain must be internal.
&
Used in the saturated $4\times64$ control.
This reconstructs a conventional full width carry chain through \code{RFLAGS}. \\
\bottomrule
\end{tabularx}
\end{table}

This distinction determines how well each field representation maps to the microcode interface.
The $5\times51$ representation leaves thirteen unused bits in each 64 bit limb, so a column sum absorbs many products before a carry has to leave it; carries are comparatively rare, and paying a third $\mu$op to move one as a value costs little.
The saturated $4\times64$ representation instead requires full width propagation between adjacent words on every column, which makes the cost of an individual carry the dominant term and suggests the cheaper two $\mu$op path should pay for its staging overhead.
We built the $4\times64$ control both ways rather than assume that, and measured all three versions back to back in one process at a pinned clock.
The \code{SETCC} dance form occupies $119$ triads and runs at $211$ cycles per multiplication.
The \code{ADC} form, with the bridge and the following \code{ADC} packed into shared triads, occupies $75$ triads and runs at $207$ cycles.
Fusing the schoolbook into the carry chain --- each multiply writing its high half straight to the register the column chain reads, and each low half consumed in the same triad that produced it --- takes the same algorithm to $67$ triads and $191$ cycles.
The control we report is that last version.
Two things in those numbers are worth stating plainly.
The two carry paths are within two percent of each other, so the \code{ADC} path's lower operation count per carry buys almost nothing here; what bought the remaining nine percent was scheduling the multiplies inside the carry chain, which is legal only because a published carry survives an intervening multiply.
And the saving did not track the triad count: removing $44$ triads between the first two versions was worth $4$ cycles, while the next $8$ were worth $16$, because those eight sat directly between each multiply and the addition consuming it.
Position on the dependency chain, not patch size, is what a triad costs.
All three versions remain far behind the $58$ cycle $5\times51$ multiplication, and also behind an ordinary \code{-O3} compilation of the same saturated algorithm in C, which runs at $135$ cycles on this machine.
It is the representation rather than the carry path that decides the comparison.
